\chapter{Synchronized Difference}\label{ch:synchronized-difference}

The Introduction stated the book's central claim: collective viewing does not require that everyone see the same thing; it requires that different experiences remain synchronized closely enough to count as one occasion. Parts~I through~V supplied the means. This chapter returns to the claim and asks what, precisely, an audience shares when its members do not share a picture, and what ``closely enough'' has turned out to mean.

\section{Layers of collectivity}

Attending a film with others involves several kinds of sharing that ordinary cinema delivers together and that can be separated. Write the collectivity of a screening as
\begin{equation}
\mathcal C=\langle P,\ T,\ W,\ A,\ D\rangle,
\end{equation}
where $P$ is physical co-presence, being in one room; $T$ is temporal synchronization, experiencing the work on one clock; $W$ is a shared world, the same characters, places, and events; $A$ is common anchoring, points at which every experience corresponds; and $D$ is the permitted divergence, how much the experiences are allowed to differ.

Ordinary cinema sets the first four as high as they go and $D$ to zero. Everyone is in the room, on the clock, in the world, and at every point in the same place, because everyone sees the same thing. Streaming reverses much of this: $P$ and $T$ collapse, since people watch alone and at different times, and $D$ may be large, since recommendation and personalization give different people different works. Private head-mounted displays can keep $P$ and even $T$ while making $D$ unlimited and $A$ absent.

Selective cinema occupies a position none of these do. It keeps $P$ and $T$ as high as ordinary cinema, keeps $W$ and $A$ high by design, and makes $D$ positive but bounded. The audience is neither a mass that sees one thing nor a scatter of individuals seeing unrelated things. It is a \emph{coordinated plurality}: people in one place, on one clock, in one story, seeing it differently.

\section{What is shared}

It is worth listing what the members of a selective audience actually share, because the list is longer than it first appears.

They share the room, its darkness, its seats, its exits. They share the clock: the film begins and ends for everyone at once, and its scenes turn together. They share the world of the story, its people and places and principal events. They share the rendezvous, the moments at which, as \Cref{ch:switch-admissibility} defined, their realizations agree on everything that will matter. Under shared or complementary sound, they share every note of the score and every sound effect, as \Cref{ch:composite-sound} described. They share whatever the realizations have in common on screen, which adaptive gating makes literally identical light. They share the composite, the view any of them sees on lifting the glasses.

And they share each other's reactions. The laughter from the comedy's viewers reaches the drama's viewers; the silence of the drama's viewers reaches the comedy's. Each viewer hears, during the film, evidence that others are having a different experience of the same event. This is not a defect to be minimized, though it must be kept from cruelty. It is the audible form of synchronized difference, and it is something no other arrangement produces. Private screens hide the difference. A common screen abolishes it. Selective cinema lets it be heard.

\section{When is it one occasion?}

The book's central claim can now be made more exact. The earlier chapters identified the conditions under which a family holds together, and they can be gathered into a definition.

\begin{definition}[Occasion-preserving family]
A variant family, as exhibited, \emph{preserves the occasion} when three conditions hold:
\begin{enumerate}
\item \emph{bounded drift}: every realization's temporal debt is cleared at declared synchronization points, so no viewer's film runs away from the common clock;
\item \emph{recurring agreement}: structural rendezvous recur, so that fact distance is periodically brought back to zero and the realizations repeatedly agree on what has happened; and
\item \emph{a common channel}: at least one element of the experience is received identically by everyone present, such as a shared score, shared footage, or a designed composite.
\end{enumerate}
\end{definition}

Each condition answers a specific way of losing the occasion. Without bounded drift, viewers are in the same room at different points in the story, and the room's shared moments no longer line up. Without recurring agreement, the realizations become different stories that happen to share characters. Without a common channel, the viewers are private audiences who share only a wall.

The definition also explains, retrospectively, several of the book's recommendations. Envelopes (\Cref{ch:genre}) and debt accounting (\Cref{ch:editing}) secure the first condition. The switch system and rendezvous (\Cref{ch:switch-admissibility}) secure the second. Complementary sound (\Cref{ch:composite-sound}) and adaptive gating secure the third. And the definition makes a prediction: a family exhibited with private sound, long unreconciled divergences, and loose timing will not feel like one screening, however good its optics.

\section{Against the feed}

Synchronized difference has an obvious modern counterpart, and the contrast is instructive. Algorithmic feeds, targeted advertising, personalized search, and recommendation systems all show different people different things, a condition that has been described as each person living inside a personalized ``filter bubble'' \cite{pariser}. They produce divergence without synchronization. Each person's experience is private, in time and in place, and the divergence is largely invisible: nobody sees the feed of the person next to them, and few know how their own was chosen.

Selective cinema produces divergence under the opposite conditions. It is co-present, synchronized, and, at least in principle, visible. The viewers know that others are seeing something else; they can hear it happening; they can discover afterward what it was. The divergence is declared rather than hidden.

That contrast is the form's critical potential. It stages, in the oldest social setting cinema has, a condition that contemporary media produce everywhere without showing it: people facing the same surface and receiving different realities. Staged openly, the condition becomes something an audience can notice, discuss, and judge. It also has the opposite potential. The same apparatus could produce divergence that is co-present but hidden, with viewers believing they saw one film. Part~VIII is concerned with keeping the first potential and preventing the second.

\section{Beyond cinema}

The structure is not confined to film. A lecture could run several explanatory streams on one screen for students at different levels, all in the room together. A museum installation could show visitors different histories of the same object. A multilingual performance could carry each language on its own stream. A civic presentation could show the same proposal from the perspective of different groups it affects, to an audience that knows each member is seeing one of several. Collaborative work around a shared display could give each participant their own layer.

Cinema remains the clearest case because it has made common exposure central to what attendance means. To go to the cinema is, conventionally, to see what everyone else in the room sees. Selective cinema separates attendance from exposure and shows that the first survives without the second, provided the synchronization is kept. Whether the same holds in lectures, galleries, and assemblies is a question for those settings. The form of the answer, bounded drift, recurring agreement, and a common channel, should carry over.
